\subsection{选址优化模型}

对候选点 \(s\) 和可能源位置 \(g\in\Omega_1\)，记两条观测射线的锐交会角为
\begin{equation}
  \phi(s,g)=\arccos\frac{|(g-S_1)^{\mathsf T}(g-s)|}
  {\|g-S_1\|_2\,\|g-s\|_2}.
\end{equation}
交会角接近 \(90\degree\) 时几何精度较好。定义归一化评分
\begin{equation}
  J(s)=w_1\min_{g\in\Omega_1}\sin\phi(s,g)
  +w_2\frac{\|s-S_1\|_2}{L_{\max}}
  -w_3\frac{\|s-S_1\|_2}{5T_{\mathrm{ref}}}
  -w_4P_{\mathrm{recv}}(s),
\end{equation}
其中 \(P_{\mathrm{recv}}\) 衡量超出保守接收域的风险。最终选择
\begin{equation}
  S_2^*=\arg\max_{s\in\mathcal{C}}J(s).
  \label{eq:p2-objective}
\end{equation}

\begin{modelalgorithm}{第二检测点选择}
  \AlgoIn{\(S_1,\theta_1,\Omega_1\) 与评分参数}
  \AlgoOut{第二检测点 \(S_2^*\)}
  \begin{enumerate}[leftmargin=2.5em]
    \AlgoStep{在目标圆内生成候选网格并按式~\eqref{eq:candidate-region} 筛选。}
    \AlgoStep{对每个候选点，在 \(\Omega_1\) 边界上采样并计算最坏交会角。}
    \AlgoStep{计算接收风险、移动成本与综合评分 \(J\)。}
    \AlgoStep{返回评分最大的可达候选点。}
  \end{enumerate}
\end{modelalgorithm}

\placeholder{用演练实验确定权重与基线区间，并说明参数选择依据。}
